%% Chapter 3, Section 3.2: Importance Sampling
%% A First Course in Monte Carlo Methods — Sanz-Alonso & Al-Ghattas

\section{Importance Sampling}
\label{sec:3.2}

Now we study how to approximate $\calI_f[h]$ using samples from a proposal
distribution $g \neq f$. There are two distinct motivations to sample from a
proposal distribution $g$, rather than from the target density $f$.

\medskip
\noindent\textbf{Motivation 1: Variance Reduction.}
If the proposal distribution is cleverly chosen, one can obtain an estimator with
a relatively lower variance. From a historical perspective, variance reduction was
the original motivation for introducing importance sampling; in particular, the
design of proposal distributions for the simulation of rare events has been
extensively researched. A natural question in this area is how best to choose the
proposal distribution in order to compute the probability of a rare event under
the target. In such a context, the test function $h$ would be the indicator of a
small-probability set.

\medskip
\noindent\textbf{Motivation 2: Intractability of Target.}
In some applications, sampling directly from $f$ may not be possible, but one
may have access to samples from a different distribution $g$. This scenario
arises for instance in sequential signal estimation, where prior samples obtained
by a stochastic dynamics model are used to compute expectations with respect to a
posterior distribution which incorporates observed data. A natural question in
this area is: given a proposal and a target, what is the worst-case error of
importance sampling over a given class of test functions?

\medskip
In its simplest form, importance sampling is based on the observation that if
$f$, $g$ and $h$ have compatible supports, then
\[
  \calI_f[h]
  = \int_E h(x)f(x)\,dx
  = \int_E h(x)\frac{f(x)}{g(x)}\,g(x)\,dx
  = \calI_g\!\left[h\frac{f}{g}\right].
\]
This identity suggests that $\calI_f[h]$ may be approximated using classical
Monte Carlo integration by sampling the density $g$ and considering $hf/g$ as
the test function.

\begin{figure}[htbp]
\centering
\includegraphics[width=0.75\textwidth]{figures/fig_03_1}
\caption{\textit{In classical Monte Carlo (left), samples from $f$ are given
equal weights. In importance sampling (right), samples from $g$ are given
weights proportional to $f/g$, represented by the radii of the dots.}}
\label{fig:3.1}
\end{figure}

\begin{algorithm}[H]
\caption{Importance Sampling}
\label{alg:3.2}
\begin{algorithmic}[1]
\Require Target distribution $f$, proposal distribution $g$, test function $h$,
         sample size $N$.
\State Sample $X^{(1)}, \ldots, X^{(N)} \iid g$.
\State Set
  \[
    w\!\left(X^{(n)}\right) := \frac{f\!\left(X^{(n)}\right)}{g\!\left(X^{(n)}\right)}.
  \]
\Ensure Estimator
  $\calI_f^{\mathrm{is}}[h] := \dfrac{1}{N}\displaystyle\sum_{n=1}^{N}
   w\!\left(X^{(n)}\right) h\!\left(X^{(n)}\right) \approx \calI_f[h]$.
\end{algorithmic}
\end{algorithm}

Note that the importance sampling estimator $\calI_f^{\mathrm{is}}[h]$ depends
on the choice of proposal $g$, but we do not include that dependence in our
notation. In addition to studying the importance sampling estimator
$\calI_f^{\mathrm{is}}[h]$, we will study an autonormalized implementation
motivated by the two following drawbacks of Algorithm~\ref{alg:3.2}.

First, note that in order to compute the weights
\[
  w\!\left(X^{(n)}\right) := \frac{f\!\left(X^{(n)}\right)}{g\!\left(X^{(n)}\right)}
\]
we need to be able to evaluate the ratio $f/g$, which may in practice involve
evaluating both the target and the proposal with their appropriate normalizing
constants.

Second, note that
\[
  \calI_g[w]
  = \int w(x)\,g(x)\,dx
  = \int \frac{f(x)}{g(x)}\,g(x)\,dx
  = 1,
\]
but in general
\[
  \frac{1}{N}\sum_{n=1}^{N} w\!\left(X^{(n)}\right) \neq 1.
\]
In other words, importance sampling does not estimate exactly the constant test
function $h \equiv 1$. These two observations motivate the use of
\emph{autonormalized} importance sampling.

\begin{algorithm}[H]
\caption{Autonormalized Importance Sampling}
\label{alg:3.3}
\begin{algorithmic}[1]
\Require Target $f$, proposal $g$, test function $h$, sample size $N$.
\State Sample $X^{(1)}, \ldots, X^{(N)} \iid g$.
\State Set
  \[
    w^*\!\left(X^{(n)}\right)
    := \frac{w\!\left(X^{(n)}\right)}{\displaystyle\sum_{l=1}^{N} w\!\left(X^{(l)}\right)}.
  \]
\Ensure Autonormalized estimator
  $\calI_f^{\mathrm{ais}}[h]
   := \displaystyle\sum_{n=1}^{N}
      w^*\!\left(X^{(n)}\right) h\!\left(X^{(n)}\right)
   \approx \calI_f[h]$.
\end{algorithmic}
\end{algorithm}

The quantities $w^*(X^{(n)})$ are known as \emph{normalized} weights because
they add up to $1$. Notice that $w$ appears in the numerator and the denominator
of $w^*$, and so $w^*$ can be evaluated as long as $w(x) = f(x)/g(x)$ is known
up to a multiplicative constant. Thus, autonormalized importance sampling can be
implemented as long as $f$ and $g$ are known up to a normalizing constant, which
greatly extends the applicability of the method. As we will see, the advantages
afforded by the autonormalized estimator come at the cost of introducing a bias.

%% -----------------------------------------------------------------------
\subsection{Error Bounds for Importance Sampling}
\label{sec:3.2.1}

In this subsection we analyze the importance sampling estimator
$\calI_f^{\mathrm{is}}[h]$. The presentation is focused on the variance
reduction motivation for importance sampling; thus we aim at understanding the
error for a given test function, and we will study the choice of proposal
distribution to reduce the variance of the Monte Carlo estimator.

The following result is parallel to Theorem~\ref{thm:3.1}. As in classical
Monte Carlo, the importance sampling estimator is unbiased and its mean squared
error is of order $1/N$. Now, however, the mean squared error depends on the
proposal distribution $g$.

\begin{theorem}[Importance Sampling Error]
\label{thm:3.2}
The following holds:
\begin{enumerate}
  \item $\Expect\!\left[\calI_f^{\mathrm{is}}[h]\right] = \calI_f[h]$;
  \item $\Var\!\left[\calI_f^{\mathrm{is}}[h]\right]
         = \dfrac{1}{N}\,\Var_{X \sim g}\!\left[h(X)\,w(X)\right]$.
\end{enumerate}
\end{theorem}

\begin{proof}
It is a corollary of Theorem~\ref{thm:3.1} since $\calI_f[h] = \calI_g[hw]$
and $\calI_f^{\mathrm{is}}[h] = \calI_g^{\mathrm{mc}}[hw]$.
\end{proof}

Theorem~\ref{thm:3.2} suggests that in order to approximate $\calI_f[h]$ for
given target $f$ and test function $h$ one should choose the proposal $g$ that
minimizes
\[
  \Var_{X \sim g}\!\left[h(X)\,w(X)\right].
\]
The next result is interesting for historical reasons and to develop intuition.
As will be discussed below, its practical significance is rather limited.

\begin{theorem}[Importance Sampling Optimal Proposal]
\label{thm:3.3}
The proposal $g$ that minimizes the variance of $\calI_f^{\mathrm{is}}[h]$ is
\[
  g^*(x) = \frac{|h(x)|\,f(x)}{\displaystyle\int_E |h(t)|\,f(t)\,dt}.
\]
\end{theorem}

\begin{proof}
We have that
\[
  N\,\Var\!\left[\calI_f^{\mathrm{is}}[h]\right]
  = \Var_g\!\left[h(X)\frac{f(X)}{g(X)}\right]
  = \Expect_g\!\left[\left(h(X)\frac{f(X)}{g(X)}\right)^2\right]
    - \Expect_g\!\left[h(X)\frac{f(X)}{g(X)}\right]^2.
\]
Since $\Expect_g\!\bigl[h(X)\tfrac{f(X)}{g(X)}\bigr]^2 = \calI_f[h]^2$ does
not depend on $g$, we only need to minimize the first term in the right-hand
side. Note that:

\begin{enumerate}[label=(\roman*)]
\item Plugging $g = g^*$ into said term gives
\begin{align*}
  \Expect_{g^*}\!\left[\left(h(X)\frac{f(X)}{g^*(X)}\right)^2\right]
  &= \int \frac{h(x)^2 f(x)^2}{g^*(x)}\,dx \\
  &= \int \frac{h(x)^2 f(x)^2}{|h(x)|\,f(x)}\,dx
     \int |h(t)|\,f(t)\,dt
   = \left(\int |h(x)|\,f(x)\,dx\right)^2.
\end{align*}

\item Applying Jensen's inequality we have that, for any $g$,
\[
  \Expect_g\!\left[\left(\frac{h(X)f(X)}{g(X)}\right)^2\right]
  \geq
  \Expect_g\!\left[\frac{|h(X)|\,f(X)}{g(X)}\right]^2
  = \left(\int |h(x)|\,f(x)\,dx\right)^2.
\]
\end{enumerate}
Combining (i) and (ii) gives the result.
\end{proof}

Theorem~\ref{thm:3.3} is of little practical relevance: in order to use $g^*$
as a proposal we need to be able to evaluate $g^*$, which involves computing an
integral similar to the one we originally wanted to compute! In fact, if $h$ is
positive, then the denominator in the definition of $g^*$ equals $\calI_f[h]$;
in that case, the proof of Theorem~\ref{thm:3.3} implies that
$\calI_f^{\mathrm{is}}[h]$ has zero variance, and a direct calculation using
the definition of $\calI_f^{\mathrm{is}}[h]$ shows that
$\calI_f^{\mathrm{is}}[h] = \calI_f[h]$. Despite the limited practical
relevance of Theorem~\ref{thm:3.3}, a useful takeaway is that importance
sampling can be super-efficient: it is possible to obtain a better (lower)
variance estimator using samples from a proposal distribution than from the
target. To further demonstrate this point, note that when $g = g^*$,
\[
  N\,\Var\!\left[\calI_f^{\mathrm{is}}[h]\right]
  = \Expect_f\!\left[|h(X)|\right]^2 - \calI_f[h]^2
  \leq \Expect_f\!\left[h(X)^2\right] - \Expect_f\!\left[h(X)\right]^2
  = \Var_f\!\left[h(X)\right]
  = N\,\Var\!\left[\calI_f^{\mathrm{mc}}[h]\right].
\]
Another useful takeaway is that samples from the proposal should concentrate on
regions where $|h|$ is large, as illustrated in the following example.

\begin{example}[Computing a Gaussian Tail Probability]
\label{ex:3.4}
Let $p = \Prob(X > 4)$ where $X \sim \Normal(0,1)$ with corresponding p.d.f.\
$f$. Using classical Monte Carlo we can approximate $p$ by
\[
  \hat{p} = \frac{1}{N}\sum_{n=1}^{N} \Ind{X^{(n)}>4},
  \qquad X^{(1)},\ldots,X^{(N)} \iid \Normal(0,1).
\]
However, $N$ would have to be very large to get an estimate that differs from
$0$. An alternative is to generate a sample of i.i.d.\ $\Exp(1)$ random
variables translated to the right by $4$, with corresponding p.d.f.\ $g$. The
weights given to the samples would be determined by the weight function
\[
  w(x) = \frac{f(x)}{g(x)}
  = \frac{1}{\sqrt{2\pi}}\exp\!\left(-\tfrac{1}{2}x^2 + (x-4)\right).
\]
Figure~\ref{fig:3.2} shows an experiment where sampling from $g$ rather than $f$
is advantageous.
\end{example}

\begin{figure}[htbp]
\centering
\includegraphics[width=0.75\textwidth]{figures/fig_03_2}
\caption{\textit{Approximation of $p = \Prob(X > 4)$ with $X \sim \Normal(0,1)$
using classical Monte Carlo and importance sampling with a shifted exponential.
The green line shows the true value, computed with Python. The approximation
with importance sampling is accurate and stable with $N = 10^4$ samples, while
classical Monte Carlo requires around $N = 10^6$ samples to reach a similar
level of accuracy.}}
\label{fig:3.2}
\end{figure}

%% -----------------------------------------------------------------------
\subsection{Error Bounds for Autonormalized Importance Sampling}
\label{sec:3.2.2}

Recall that autonormalized importance sampling can be used as long as the ratio
$f/g$ can be evaluated up to a multiplicative constant. Here, we analyze
autonormalized importance sampling in the following setting:
\begin{itemize}
  \item Target $f(x) = \tilde{f}(x)/Z$, where $f$ is a normalized distribution
        and $\tilde{f}$ is unnormalized.
  \item Normalized proposal distribution $g(x)$.
\end{itemize}
Our analysis is motivated by applications where sampling $f$ directly may not be
possible. We are interested in estimating $\calI_f[h]$ for many different test
functions $h$ or for a test function $h$ that is unknown before designing the
algorithm. Thus, we will be concerned with bounding the \emph{worst-case error}
over a \emph{family} of test functions. We will work with the family of bounded
functions, but similar results can be obtained for other classes of test
functions.

Autonormalized importance sampling is based on writing $\calI_f[h]$ as a ratio
of expectations with respect to $g$ and approximating both expectations with
classical Monte Carlo. Precisely,
\[
  \calI_f[h]
  = \frac{\calI_g[h(X)\,\tilde{w}(X)]}{\calI_g[\tilde{w}(X)]}
  \approx
  \frac{\calI_g^{\mathrm{mc}}[h\tilde{w}]}{\calI_g^{\mathrm{mc}}[\tilde{w}]}
  = \calI_f^{\mathrm{ais}}[h],
\]
where $\tilde{w}(x) = \tilde{f}(x)/g(x)$. Note that the denominator
$\calI_g^{\mathrm{mc}}[\tilde{w}]$ is an approximation of the unknown
normalizing constant $Z$ of the density $f$. The next result is parallel to
Theorems~\ref{thm:3.1} and~\ref{thm:3.2}. Note that, in contrast to classical
Monte Carlo and importance sampling, autonormalized importance sampling gives a
\emph{biased} estimator, since the ratio of two unbiased estimators is in
general biased.

\begin{theorem}[Autonormalized Importance Sampling Error]
\label{thm:3.5}
Let $h : E \to \R$ with $|h|_\infty := \sup_{x \in E} |h(x)| \leq 1$. The
following holds:
\begin{enumerate}
  \item $\displaystyle\left|\Expect\!\left[\calI_f^{\mathrm{ais}}[h]
        - \calI_f[h]\right]\right|
        \leq \frac{2}{N}\,
        \frac{\calI_g[\tilde{w}^2]}{\calI_g[\tilde{w}]^2}$;
  \item $\displaystyle\Expect\!\left[\left(\calI_f^{\mathrm{ais}}[h]
        - \calI_f[h]\right)^2\right]
        \leq \frac{4}{N}\,
        \frac{\calI_g[\tilde{w}^2]}{\calI_g[\tilde{w}]^2}$.
\end{enumerate}
\end{theorem}

\begin{proof}
Note that
\begin{align*}
  \calI_f^{\mathrm{ais}}[h] - \calI_f[h]
  &= \calI_f^{\mathrm{ais}}[h]
     - \frac{\calI_g[\tilde{w}h]}{\calI_g[\tilde{w}]} \\
  &= \left(\calI_g[\tilde{w}] - \calI_g^{\mathrm{mc}}[\tilde{w}]\right)
     \frac{\calI_f^{\mathrm{ais}}[h]}{\calI_g[\tilde{w}]}
     - \frac{\calI_g[\tilde{w}h] - \calI_g^{\mathrm{mc}}[\tilde{w}h]}%
            {\calI_g[\tilde{w}]}.
\end{align*}
This identity will be used to show the bounds for the bias and the mean squared
error. We start with the latter:
\begin{align*}
  &\Expect\!\left[\left(\calI_f^{\mathrm{ais}}[h] - \calI_f[h]\right)^2\right] \\
  &\quad\leq
  \frac{2}{\calI_g[\tilde{w}]^2}\!\left(
    \Expect\!\left[\left(\calI_f^{\mathrm{ais}}[h]\right)^2
      \left(\calI_g[\tilde{w}] - \calI_g^{\mathrm{mc}}[\tilde{w}]\right)^2\right]
    + \Expect\!\left[\left(\calI_g[\tilde{w}h] - \calI_g^{\mathrm{mc}}[\tilde{w}h]\right)^2\right]
  \right) \\
  &\quad\leq
  \frac{2}{\calI_g[\tilde{w}]^2 N}
  \left(\Var_g[\tilde{w}] + \Var_g[\tilde{w}h]\right) \\
  &\quad\leq
  \frac{4}{N}\,\frac{\calI_g[\tilde{w}^2]}{\calI_g[\tilde{w}]^2}.
\end{align*}
We have used (i) for any $a, b \in \R$, $(a-b)^2 \leq 2(a^2+b^2)$; (ii)
Theorem~\ref{thm:3.1} for mean squared error of classical Monte Carlo; (iii)
that $h$ is bounded by $1$; and (iv) the bound
$\Var[\tilde{w}] \leq \Expect_g[\tilde{w}^2] = \calI_g[\tilde{w}^2]$.

Now we prove the result for the bias:
\begin{align*}
  \left|\Expect\!\left[\calI_f^{\mathrm{ais}}[h] - \calI_f[h]\right]\right|
  &= \frac{1}{\calI_g[\tilde{w}]}
     \left|\Expect\!\left[
       \left(\calI_f^{\mathrm{ais}}[h] - \calI_f[h]\right)
       \left(\calI_g[\tilde{w}] - \calI_g^{\mathrm{mc}}[\tilde{w}]\right)
     \right]\right| \\
  &\leq
  \frac{1}{\calI_g[\tilde{w}]}
  \left(\Expect\!\left[\left(\calI_f^{\mathrm{ais}}[h] - \calI_f[h]\right)^2\right]\right)^{1/2}
  \!\left(\Expect\!\left[\left(\calI_g[\tilde{w}] - \calI_g^{\mathrm{mc}}[\tilde{w}]\right)^2\right]\right)^{1/2} \\
  &\leq
  \frac{1}{\calI_g[\tilde{w}]}
  \left(\frac{4}{N}\,\frac{\calI_g[\tilde{w}^2]}{\calI_g[\tilde{w}]^2}\right)^{1/2}
  \!\left(\frac{\calI_g[\tilde{w}^2]}{N}\right)^{1/2} \\
  &\leq
  \frac{2}{N}\,\frac{\calI_g[\tilde{w}^2]}{\calI_g[\tilde{w}]^2}.
\end{align*}
For the first equality we used that
$\Expect\!\bigl[\calI_g[\tilde{w}] - \calI_g^{\mathrm{mc}}[\tilde{w}]\bigr] = 0$,
for the first inequality Cauchy--Schwarz, and for the second inequality the bound
for the mean squared error.
\end{proof}

\begin{remark}
\label{rem:3.6}
The autonormalized importance sampling estimator is biased, but its bias can be
uniformly bounded over the class of bounded test functions $h : E \to \R$ and it
is of order $\tfrac{1}{N}$. The mean squared error of the autonormalized
importance sampling estimator is of order $\tfrac{1}{N}$ as well. Hence, since
the mean squared error equals the sum of the squared bias and the variance,
i.e.\
\[
  \mathrm{MSE} = \mathrm{bias}^2 + \mathrm{variance},
\]
it follows that for large $N$ the mean squared error is dominated by the
variance term.
\end{remark}

\medskip
\noindent\textbf{Autonormalized Importance Sampling and Closeness Between Target
and Proposal.}
The upper bounds on the bias and mean squared error depend on
\[
  \zeta := \frac{\calI_g[\tilde{w}^2]}{\calI_g[\tilde{w}]^2}.
\]
The value of $\zeta$ quantifies the variability on the weights. It holds that
$\zeta \geq 1$ with equality only if $\tilde{w}$ is constant, i.e.\ if $f = g$.
Larger values of $\zeta$ imply worse behavior of importance sampling over the
class of bounded test functions. Thus, if no knowledge of the relationship
between the test function $h$ and the target is available, it is important to
choose a proposal that is close to the target, so that $\zeta$ is small.

The quantity $\zeta$ is intimately related to the $\chi^2$-divergence between
the target and the proposal. The $\chi^2$-divergence is a way to quantify the
closeness between two densities, defined by
\[
  d_{\chi^2}(f\|g)
  := \int_E \left(\frac{f(x)}{g(x)} - 1\right)^2 g(x)\,dx.
\]
Therefore, we have that $d_{\chi^2}(f\|g) = \zeta - 1$. Hence,
Theorem~\ref{thm:3.5} implies that for autonormalized importance sampling to
perform accurately over the class of bounded test functions, the
$\chi^2$-divergence between the target and the proposal needs to be
small/moderate.

\medskip
\noindent\textbf{Effective Sample Size.}
The effective sample size
\[
  \mathrm{ESS}
  := \frac{1}{\displaystyle\sum_{n=1}^{N} w^*\!\left(X^{(n)}\right)^2},
  \qquad
  w^*\!\left(X^{(n)}\right)
  := \frac{\tilde{w}\!\left(X^{(n)}\right)}%
          {\displaystyle\sum_{n=1}^{N} \tilde{w}\!\left(X^{(n)}\right)}
\]
is widely used by practitioners as a diagnostic for the performance of
(autonormalized) importance sampling. Note that
\begin{itemize}
  \item $\mathrm{ESS} = 1$ if there is $n \in \{1,\ldots,N\}$ such that
        $w^*(X^{(n)}) = 1$ (in which case all other normalized weights are
        zero).
  \item $\mathrm{ESS} = N$ if $w^*(X^{(n)}) = \tfrac{1}{N}$ for all
        $n \in \{1,\ldots,N\}$.
  \item $\mathrm{ESS} \in [1, N]$ for any configuration of weights.
\end{itemize}
The effective sample size $\mathrm{ESS}$ is intimately related to $\zeta$ and
$d_{\chi^2}(f\|g)$. Indeed,
\[
  \frac{\mathrm{ESS}}{N}
  = \frac{1}{N\sum w^*(X^{(n)})^2}
  = \frac{\bigl(\sum \tilde{w}(X^{(n)})\bigr)^2}{N\sum \tilde{w}(X^{(n)})^2}
  = \frac{\left(\dfrac{\sum \tilde{w}(X^{(n)})}{N}\right)^2}%
         {\dfrac{\sum \tilde{w}(X^{(n)})^2}{N}}
  \approx
  \frac{\calI_g^{\mathrm{mc}}[\tilde{w}]^2}{\calI_g^{\mathrm{mc}}[\tilde{w}^2]},
\]
and so
\[
  \frac{N}{\mathrm{ESS}}
  = \frac{\calI_g^{\mathrm{mc}}[\tilde{w}^2]}{\calI_g^{\mathrm{mc}}[\tilde{w}]^2}
  \approx \zeta.
\]
It is important to note, however, that $\mathrm{ESS}$ does not include any test
function $h$ in its definition. Therefore, the performance of (autonormalized)
importance sampling with any given test function cannot be fully understood via
the effective sample size.

\medskip
\noindent\textbf{Autonormalized Importance Sampling and Curse of Dimension.}
Consider (as for rejection sampling) the i.i.d.\ setting
\[
  \tilde{f}_d(x)
  = \tilde{f}(x_1) \times \cdots \times \tilde{f}(x_d),
  \quad \tilde{f}: \R \to \R,
  \quad x = (x_1,\ldots,x_d) \in \R^d,
\]
\[
  g_d(x) = g(x_1) \times \cdots \times g(x_d),
  \quad g: \R \to \R,
\]
where the tildes stand for unnormalized densities. Then,
\[
  \frac{\calI_{g_d}[\tilde{w}_d^2]}{\calI_{g_d}[\tilde{w}_d]^2}
  = \left(\frac{\calI_g[\tilde{w}^2]}{\calI_g[\tilde{w}]^2}\right)^d.
\]
This identity along with Theorem~\ref{thm:3.5} suggest that in order to have the
same accuracy in importance sampling with proposal $g_d$ and target $f_d$ across
a sequence of problems with increasing $d$, one needs to increase the number of
samples exponentially with $d$.

\medskip
\noindent\textbf{Pros and Cons.}
As opposed to classical Monte Carlo, (autonormalized) importance sampling can be
implemented when $f$ cannot be sampled from and, in contrast to rejection
sampling, no samples are thrown away. Importance sampling can have smaller
variance than classical Monte Carlo integration, but finding an adequate proposal
that gives variance reduction can be challenging. Autonormalized importance
sampling can have smaller variance than both importance sampling and classical
Monte Carlo, but is biased. On the other hand, importance sampling requires
evaluating the normalized target and proposal distributions, but is unbiased. If
the target and the proposal are far apart (which is typically the case in
high-dimensional problems) the weights have a large variance. Then, the effective
sample size will be small, and there will be test functions for which importance
sampling performs poorly.
